% vim: tw=78 ai spell spelllang=en

\documentclass[10pt]{llncs}
\usepackage{color}
\usepackage{url}
\usepackage{xspace}
\usepackage{graphicx}
\usepackage{listings}

\newcommand\Color[2]{\textcolor{#1}{#2}}
\newcommand{\nb}[1]{\marginpar{\scriptsize #1}}
\renewcommand\note[1]{\nb{\Color{blue}{#1}}}
\newcommand\todo[1]{\nb{\Color{red}{#1}}}

\newcommand{\COC}{\textsc{Coccinelle}\xspace}
\newcommand{\COV}{\textsc{Coverity}\xspace}
\newcommand{\SM}{\textsc{Smatch}\xspace}
\newcommand{\SP}{\textsc{Sparse}\xspace}
\newcommand{\ST}{\textsc{Stanse}\xspace}

\lstset{tabsize=2,numbers=none,showstringspaces=false,breaklines=true,
        breakatwhitespace=true,escapeinside={(*@}{@*)},
	basicstyle=\scriptsize}

\begin{document}

\title{Kernel Checking: a Case Study}
\author{Ji\v{r}\'\i~Slab\'y and Jan~Strej\v{c}ek}
\institute{Faculty of Informatics, Masaryk University, Brno, Czech Republic
\email{\{slaby,strejcek\}@fi.muni.cz}}

\maketitle

\begin{abstract}
We present a case study of static analysis tools run on the Linux kernel.
There are totally of XXX tools we introduce, describe in more detail and
evaluate on the kernel. Also overall statistics are summed in the results
section.
\end{abstract}

\section{Introduction}
\label{sec:introduction}
In the ideal world, every software product is perfectly working. This can be
accomplished only without presence of bugs. To have a code containing no
errors is of great concern since middle of 60's\note{fakt?}. 

There is a variety of techniques allowing for pruning bugs out of code. Most
of the invented algorithms is implemented and provided to be used. Some are
freely available or even open-source like \SM~\cite{smatch},
\SP~\cite{sparse}, \ST~\cite{stanse} and \textsc{FindBugs}~\cite{findbugs}.
Others check open-source code, but are not available to programmers. The most
known tool belonging to this category is \COV~\cite{coverity}. Finally, the
tools which are commercial only include \textsc{Klocwork}~\cite{klocwork} and
\textsc{Codesonar}\cite{codesonar}.

Absence of bugs is very important in all projects. However a crash of instant
messaging client or web browser is surely less important than a crash causing
a banking system to collapse. Unfortunately the mission critical software like
the one in banks is very complex and not all tools may be used for finding
bugs there.

Banking systems are not the only at the same level of complexity and severity.
There are many other projects, air navigation systems, nuclear power plant
controlling systems or kernels of operating systems. In this work, we present
a case study of bug-finding tools run on the Linux kernel.

Nowadays many bug-finding tools may be and are run on the Linux Kernel.

In the section~\ref{sec:tools} tools which might be used on the Linux kernel
are described. The next section~\ref{sec:results} contains comparison and
evaluation.

\section{Tools}
\label{sec:tools}

This section contains a brief description of tools we evaluate on the Linux
kernel. We also present types of bugs which each tool can catch.

\subsection{Coverity}
It is a product originally developed at the XXX university. It is based solely
on static analysis.

\subsection{Smatch and Sparse}
\SP is a simple parser with some basic checks on the top of that.  It was
developed by Linus Torvalds, the Linux kernel founder, to check some simple
properties in the kernel.

It can catch errors such invalid dereferences of addresses which were intended
to be dereferenced only via special helpers. These helpers usually enforces
the compiler to not optimize the accesses to the memory or devices themselves.
It also looks for variable definition which should be local to files but are
not marked by static modifier.

\SM is a follower of \SP. It uses its compiler, but have more eligible
techniques implemented. It is useful in finding bugs in memory accesses like
NULL pointer dereference or memory leaks.

It is very easy to add a "checker" into \SM. It is a matter of adding a hook
and specifying what to do (a function) and when (on fucntion entry/return,
dereference\note{fakt?}, etc.).

Both of the tools are intra-procedural only and loops are unrolled exactly
3\note{fakt?} times. Internal structures cannot be by default used by the
checkers, only hooks may be used.

\subsection{Stanse}
\ST was developed at the Masaryk University. It is rather a framework intended
for implementation of checking algorithms. Currently there are three checkers
which might be used on the Linux kernel:
\begin{itemize}
\item \textsc{AutomatonChecker} -- pairing of locking functions, memory 
\item \textsc{ThreadChecker} -- finds potential deadlocks in the code
\item \textsc{ReachabilityChecker} -- simple checker to find unreachable code
\end{itemize}

\ST analyses the code inter-procedurally, provides checker internal
representations of the code.

\subsection{Coccinelle}
There exist also tools developed for purposes other than finding bugs which
were later reinvented for bug-finding. To this set we do not count simple
pattern finding languages like \textsc{grep} or \textsc{awk}. We rather mean
\COC~\cite{coccinelle} and similar.

Originally, \COC was a semantic patcher. It was an instrument to automatically
change some patterns to some other patterns in the code.  Renaming a function
is a good example of use.

Nowadays, the tool is used also for finding simple bugs in the kernel code. It
can catch missing unlocking functions or dereferences after the memory is
freed.

\subsection{Other Tools}
There are probably many other tools capable of parsing the Linux kernel and we
miss in this study. It is mostly of much energy invested into the initial
setup. All of the tools we present here are easy to set up (without much
burder) and even easier to be run.

\section{Evaluation}
\label{sec:results}
\COV's results are taken over from~\cite{billion}. \COC, \SM, \SP and \ST was
evaluated on the Linux kernel version 2.6.XXX. In each of the following
sections we describe step-by-step how the analysis was performed.

\subsection{Setup}

\subsection{Run}

\subsection{Results}

\section{Conclusion}
\label{sec:conclusion}
In this paper we described some static analysis tools which are used to
finding bugs in the Linux kernel. We describe the tools slightly and then we
provide results of the Linux kernel runs. These include description of kinds
of bugs found and statistics about the bugs themselves.

It turns out that the only commercial tools we had the results from can find
variety of errors with minimal effort spent on the crawling thorough the bug
reports thanks to low false positive rate.

\subsubsection{Acknowledgements}
\label{sec:acknowledgements}

All authors are supported by the research centre Institute for Theoretical
Computer Science (ITI), project No. 1M0545.

\bibliographystyle{plain}
\bibliography{kernel_checking}

\end{document}
